\section{Indexing and File Structures}

\subsection{Index Implementation}
To optimize query performance, several indexes were created. The choices include B+ Tree indexes for range queries and Hash indexes for exact lookups where applicable.

\lstinputlisting[language=SQL, caption={Index Creation (\texttt{sql/indexing.sql})}]{../sql/indexing.sql}

\subsection{Performance Analysis}
This section discusses the performance impact of the implemented indexes.

\subsubsection{Indexed vs. Non-Indexed Query Performance}
We compared the execution time of several key queries before and after creating indexes. For example, finding all appointments for a specific patient, retrieving a doctor's appointments within a date range, or listing a patient's lab history sorted by date.

\textbf{Example Query: Find all appointments for a specific patient.}
\begin{itemize}
    \item \textbf{Without Index:} The database may perform a sequential scan on the \texttt{appointment} table. See the \texttt{EXPLAIN ANALYZE} output for the plan and timing.
    \item \textbf{With B+ Tree Index on \texttt{appointment(patient\_id)}:} The database can use an index scan to quickly locate matching rows. See the \texttt{EXPLAIN ANALYZE} output for the plan and timing.
\end{itemize}

On larger datasets, indexes are expected to improve query performance by reducing the amount of data scanned and by avoiding expensive sorts. With a small sample dataset, PostgreSQL may still choose sequential scans, which is not an error but a cost-based planner decision.
